\documentclass[10pt,a4paper]{amsart}
\usepackage[margin=19mm]{geometry}
\usepackage{amsmath,amssymb,mathtools}
\usepackage[scr=rsfs,cal=euler]{mathalfa}
\usepackage{xfrac}
\usepackage[unicode,hidelinks,bookmarks=false]{hyperref}
\newcommand{\ie}{i.e.\ }
\newcommand{\cf}{cf.\ }
\newcommand{\resp}{respectively\ }
\newcommand{\Subsection}{\subsection}
\providecommand{\nobreakdash}{\nobreak\hbox{-}}
\setlength{\emergencystretch}{2em}
\multlinegap0pt
\usepackage{bm}
\usepackage{bbm}
\usepackage{dsfont}
\usepackage{xspace}
\usepackage{cancel}
\usepackage{mathtools}
\usepackage{amsthm}
\newtheorem{theorem}{Theorem}
\newtheorem{corollary}[theorem]{Corollary}
\newtheorem{lemma}[theorem]{Lemma}
\title[A structural proof of the Karpelevi\v{c} theorem]{A structural proof of the Karpelevi\v{c} theorem}
\author{Brecht Verbeken, Vincent Ginis}
\date{Source: arXiv:2609.26058v2 (CC BY 4.0)}
\begin{document}
\maketitle

\noindent\emph{Source and licence.} A structural proof of the Karpelevi\v{c} theorem. Version arXiv:2609.26058v2 (CC BY 4.0), distributed under the Creative Commons Attribution 4.0 International licence, as recorded in the arXiv licence metadata for this exact version and shown on the abstract page https://arxiv.org/abs/2609.26058v2. Full rights notice: licenses/arxiv-2609.26058-CC-BY-4.0.txt.\par\smallskip

\noindent\textit{Editorial scope.} Counted unit: the Corollary whose source label is \texttt{lem:positive-mediant} in the article (source0.tex lines 1674--1678, source label \texttt{lem:positive-mediant}), delivered together with the article's own complete original proof of it (source0.tex lines 1679--1778, one \texttt{proof} environment of 100 lines). No mathematics is written, repaired or re-derived: statement and proof are the article's own text line for line. Required context retained uncounted: lem:add-factor (lines 1624--1629); thm:sector-difference (lines 1648--1657). Every definition, display and label of those lines is retained unchanged. Omitted from this package and not redistributed: 1--1623, 1630--1647, 1658--1673, 1779--2644. Nothing inside a retained excerpt is omitted or altered: every retained body is delivered exactly as the article's lines are written, so no omission receipt of any kind accompanies this package. Citations: the retained excerpts carry no citation command at all, so no bibliography entry of the article is transcribed and no reference is redistributed. Reference adaptation: none was needed and none was made; every reference inside the delivered unit resolves inside it, so no unresolved reference or citation is produced. Printed slips kept literal: no printed word, symbol, hypothesis or number of the counted statement or of its proof was altered. Layout and class: the article's own class is \texttt{article} with packages including fontenc, lmodern, textcomp, amsmath, amssymb, amsthm, mathtools, geometry, microtype, xcolor, tikz, flafter, needspace, hyperref; the wrapper is the lane's pinned \texttt{amsart} editorial wrapper at 10pt on A4, which restates the theorem-like environments the retained text needs and adds the article's own macro definitions that the retained text needs (none), with extra wrapper packages (bm, bbm, dsfont, xspace, cancel, mathtools). The delivered numbering is the unit's own. Assets: no retained span contains a figure environment, a graphics-inclusion command, a PDF-page inclusion, a table or a TikZ picture; no image file is copied into this package, the package declares no asset dependency and no image file is redistributed with it. Licence: version arXiv:2609.26058v2 of this article is distributed under the Creative Commons Attribution 4.0 International licence, as recorded in the arXiv licence metadata for this exact version and shown on the abstract page https://arxiv.org/abs/2609.26058v2. A full rights notice is delivered at licenses/arxiv-2609.26058-CC-BY-4.0.txt. Characters: every retained excerpt is pure ASCII, so the delivered engine, XeLaTeX with its default Latin Modern text font, reports no missing character.
\begin{lemma}[Adding a factor]\label{lem:add-factor}
Fix $p/q<y<r/s$, $rq-ps=1$, $q<s$.
For an integer $m\ge1$, suppose
$A=2\pi(qy-p)>0$, $B_m=2\pi(r-sy)/m>0$, and $A+B_m<\pi$.
The scalar root for count $m+1$ is strictly larger than the one for count $m$.
\end{lemma}

\begin{lemma}[A segment comparison]
\label{thm:sector-difference}
Let $f$ be holomorphic near the segment $[a,b]\subset\mathbb C$, with $a\ne b$,
and let $\chi\in\mathbb R$.  If
$\operatorname{Im}(e^{i\chi}f'(z))>0$ on $[a,b]$, then
\[
 \operatorname{Im}\!\left(e^{i\chi}
          (\bar a-\bar b)(f(a)-f(b))\right)>0.
\]
\end{lemma}

\begin{corollary}[Strict mediant refinement]\label{lem:positive-mediant}
Suppose $a/b<c/d$ are consecutive in $F_{n-1}\cap[0,1/2]$ and
$b+d=n\ge4$.  On either new open interval formed by $(a+c)/n$, the new scalar
root is strictly larger than the old one.
\end{corollary}

\begin{proof}
Reflection permits $b<d$.  Put
\[
 m=\lfloor(n-1)/b\rfloor,\qquad
 A=2\pi(bx-a),\qquad B=2\pi(c-dx)/m.
\]
The old scalar root $\rho\in(0,1)$ satisfies
\begin{equation}\label{eq:old-mediant-scalar}
 \rho^{d/m}\sin A+\rho^b\sin B=\sin(A+B),
 \qquad A,B>0,\quad A+B<\pi.
\end{equation}
The mediant splits the interval according to the sign of
\[
 mB-A=2\pi\bigl((a+c)-(b+d)x\bigr).
\]
In the left child $A<mB$, keep count $m$ initially; the smaller denominator
remains $b$, and its two angles are $A,B-A/m$.
In the right child $A>mB$, reflection puts the smaller denominator $d$
first.  Its denominators are $d,b+d$, its count is one because $b<d$,
and its angles are $mB,A-mB$.  The angle sums are respectively
$A+B-A/m<\pi$ and $A<\pi$, so both scalar equations have unique roots.
Their deficits at the old radius are
\begin{align*}
 D_L&=\sin(A+B-A/m)-\rho^{(b+d)/m}\sin A
                            -\rho^b\sin(B-A/m),\\
 D_R&=\sin A-\rho^{b+d}\sin(mB)-\rho^d\sin(A-mB).
\end{align*}
Proving these deficits positive is enough: each new scalar left side
increases strictly with the radius. To find their signs, we use the old
equation to express each deficit using powers of a single radial variable. The left
child combines exponents in units of $1/m$, whereas the right child
has count one; this dictates the substitutions
\[
 \eta=A/m,\qquad u=\rho^{b/m},\qquad v=\rho^{d/m}.
\]
Set $G_m(a,b)=(\bar a-\bar b)(a^m-b^m)$.
Then the old scalar equation reads
$v\sin A+u^m\sin B=\sin(A+B)$.
In $D_L$, the mixed power is $\rho^{(b+d)/m}=uv$.
Multiplying the old equation by $u$ replaces $uv\sin A$ by
$u\sin(A+B)-u^{m+1}\sin B$. This eliminates $v$ and gives
\[
 D_L=\sin\bigl(B+(m-1)\eta\bigr)-u\sin(B+m\eta)
       -u^m\sin(B-\eta)+u^{m+1}\sin B.
\]
These four terms are the imaginary part of
$e^{iB}(e^{-i\eta}-u)(e^{im\eta}-u^m)$.
For $D_R$, both radial terms contain $v^m$, since
$\rho^{b+d}=u^mv^m$ and $\rho^d=v^m$. We therefore eliminate $u^m$
instead, using
$u^m\sin B=\sin(A+B)-v\sin A$.
After multiplication by the positive quantity $\sin B/\sin A$, the identity
\[
 \sin(A+B)\sin(mB)+\sin B\sin(A-mB)
       =\sin A\sin\bigl((m+1)B\bigr)
\]
then gives
\[
 \frac{\sin B}{\sin A}D_R
   =\sin B-v^m\sin\bigl((m+1)B\bigr)+v^{m+1}\sin(mB).
\]
This is the imaginary part of
$e^{iB}(1-ve^{-iB})(1-v^me^{imB})$; the remaining term $-v$ is real.
Thus the two deficits are instances of the same divided difference:
\begin{equation}\label{eq:mediant-differences}
 \begin{split}
 D_L&=\operatorname{Im}\!\left(e^{iB}G_m(e^{i\eta},u)\right),\\
 D_R&=\frac{\sin A}{\sin B}\,
       \operatorname{Im}\!\left(e^{iB}G_m(1,ve^{iB})\right).
 \end{split}
\end{equation}
The factors $\bar a-\bar b$ in these expressions record the directions
of the two segments. Lemma~\ref{thm:sector-difference} tests their signs
by averaging the derivatives of $f(z)=z^m$, rotated through $B$.
The left segment joins the positive real number $u$ to $e^{i\eta}$;
the right joins $1$ to $ve^{iB}$.  Their arguments lie respectively in
$[0,\eta]$ and $[0,B]$.  Since $\eta,B<\pi$, each segment lies in an open
half-plane through the origin and avoids zero.
For the left segment, the rotated derivative $e^{iB}mz^{m-1}$ has arguments in
\[
 [B,B+(m-1)\eta]\subset(0,\pi),
 \qquad B+(m-1)\eta<A+B<\pi.
\]
For the right segment its arguments lie in $[B,mB]\subset(0,\pi)$,
because the right-child condition is $mB<A<\pi$.
These statements also cover $m=1$, when the derivative is constant.
The lemma makes both imaginary parts positive; the factor $\sin A/\sin B$
is positive as well.  Each child's scalar left side is therefore below its
target at $\rho$, so its strictly increasing dependence on radius places
the child root above $\rho$.

The right child's count is already correct.  On the left, the actual count
is $\lfloor(b+d)/b\rfloor$, whereas the comparison used
$m=\lfloor(b+d-1)/b\rfloor$.
They differ precisely when $b\mid d$; coprimality of the two denominators
makes this equivalent to $b=1$.
In that case $m=d$ and the actual count is $d+1$.
Lemma~\ref{lem:add-factor} increases the radius once more, starting with
the valid angle sum $A+B-A/m<\pi$.
\end{proof}

\end{document}
